\documentclass[conference]{IEEEtran}

\usepackage[utf8]{inputenc}
\usepackage[T1]{fontenc}
\usepackage{amsmath,amssymb}
\usepackage{graphicx}
\usepackage{booktabs}
\usepackage{siunitx}
\usepackage{array}

\graphicspath{{figuras/}}

\newcommand{\imgplaceholder}[2][0.95\columnwidth]{%
    \IfFileExists{#2}{%
        \includegraphics[width=#1]{#2}%
    }{%
        \fbox{\parbox[c][4.0cm][c]{0.90\columnwidth}{\centering
        Ejecutar \texttt{run\_all} en MATLAB para generar\\
        \texttt{#2}}}%
    }%
}

\title{Modelado y Control del Nivel de Agua en una C\'amara de Esclusa}

\author{
\IEEEauthorblockN{Nombre 1, Nombre 2}
\IEEEauthorblockA{
Universidad de los Andes, Bogot\'a, Colombia\\
correo1@uniandes.edu.co, correo2@uniandes.edu.co}
}

\begin{document}

\maketitle

\begin{abstract}
Se desarrolla un modelo y un esquema de control para regular el nivel de agua en una c\'amara de esclusa. Se plantea un modelo no lineal con p\'erdidas turbulentas, dependencia de la carga hidr\'aulica con la altura y saturaci\'on del actuador de la v\'alvula. A partir de su linealizaci\'on alrededor de un punto de operaci\'on se obtiene una planta SISO de tercer orden, con la que se dise\~na un controlador por modelo interno (IMC). Se comparan las respuestas del modelo lineal y del modelo realista en lazo abierto y lazo cerrado.
\end{abstract}

\section{Introducci\'on}
El control del nivel de agua en una esclusa es relevante porque una variaci\'on excesiva del caudal puede comprometer la seguridad de la operaci\'on y la eficiencia del proceso de llenado o vaciado. Desde el punto de vista de control, el sistema es atractivo porque combina almacenamiento, din\'amica del actuador y efectos no lineales hidr\'aulicos. En este proyecto se toma como entrada la apertura normalizada de una v\'alvula y como salida la desviaci\'on del nivel de agua respecto a un nivel nominal.

\section{Modelo no lineal}
Se definen los estados $h(t)$, $q(t)$ y $x_v(t)$ como la desviaci\'on del nivel, el caudal neto en el conducto y la apertura efectiva de la v\'alvula, respectivamente. El modelo realista propuesto es
\begin{equation}
A_c \dot{h}=q,
\end{equation}
\begin{equation}
L_h \dot{q}=K_v x_v - R_h q - R_t q|q|
- K_{h,nl}\left(\sqrt{h_0+h}-\sqrt{h_0}\right) + d(t),
\end{equation}
\begin{equation}
\tau_v \dot{x}_v + x_v = \mathrm{sat}(u).
\end{equation}
El t\'ermino $q|q|$ representa p\'erdidas turbulentas, mientras que el t\'ermino con ra\'iz modela la dependencia de la carga hidr\'aulica con la altura del agua. La saturaci\'on limita el mando del actuador.

\section{Modelo lineal aproximado}
Linealizando alrededor del equilibrio $h^\star=0$, $q^\star=0$, $x_v^\star=0$ y $d^\star=0$, se obtiene
\begin{equation}
A_c \dot{h}=q,
\end{equation}
\begin{equation}
L_h \dot{q}=K_v x_v - R_h q - K_h h + d(t),
\end{equation}
\begin{equation}
\tau_v \dot{x}_v + x_v = u.
\end{equation}
La funci\'on de transferencia entre la entrada de control y el nivel es
\begin{equation}
G(s)=\frac{H(s)}{U(s)}=
\frac{K_v}{(\tau_v s+1)(A_cL_h s^2 + A_cR_h s + K_h)}.
\end{equation}
Con $A_c=\SI{40}{m^2}$, $L_h=3$, $R_h=0.65$, $K_h=1$, $K_v=0.70$ y $\tau_v=\SI{2}{s}$, resulta
\begin{equation}
G(s)=\frac{0.70}{240s^3 + 172s^2 + 28s + 1}.
\end{equation}
Los polos son $-0.5$, $-0.1667$ y $-0.05$, de modo que la planta aproximada es estable en lazo abierto.

\section{Dise\~no del controlador}
Se adopta control por modelo interno con el filtro
\begin{equation}
F(s)=\frac{1}{(\lambda s + 1)^3}, \qquad \lambda=\SI{5}{s}.
\end{equation}
Como la planta nominal es estable e invertible, el controlador interno es $Q(s)=G^{-1}(s)F(s)$ y el controlador equivalente en realimentaci\'on queda
\begin{equation}
G_c(s)=
\frac{(\tau_v s+1)(A_cL_h s^2 + A_cR_h s + K_h)}
{K_v \lambda s(\lambda^2 s^2 + 3\lambda s + 3)}.
\end{equation}
En la planta nominal, la transferencia deseada de lazo cerrado es
\begin{equation}
T(s)=\frac{1}{(\lambda s + 1)^3},
\end{equation}
por lo que se espera una respuesta suave y sin oscilaciones fuertes.

\section{Escenarios de simulaci\'on}
Se consideran dos pruebas:
\begin{enumerate}
\item Lazo abierto con un escal\'on de amplitud $u=0.25$.
\item Lazo cerrado con la referencia por tramos $0 \rightarrow 0.25 \rightarrow -0.15 \rightarrow 0.10$ m y una perturbaci\'on sinusoidal entre $t=\SI{80}{s}$ y $t=\SI{180}{s}$.
\end{enumerate}
La perturbaci\'on representa una variaci\'on hidr\'aulica externa, por ejemplo cambios en el nivel del canal o una excitaci\'on inducida por el tr\'ansito de una embarcaci\'on.

\section{Resultados}
La Fig.~\ref{fig:pzmap} muestra el mapa de polos y ceros de la planta aproximada. La Fig.~\ref{fig:ol} compara el comportamiento lineal y no lineal en lazo abierto. La Fig.~\ref{fig:cl} presenta el seguimiento de referencia y el rechazo a perturbaciones en lazo cerrado, mientras que la Fig.~\ref{fig:u} muestra el esfuerzo de control.

\begin{figure}[!t]
\centering
\imgplaceholder{pzmap_planta_aproximada.png}
\caption{Mapa de polos y ceros de la planta aproximada.}
\label{fig:pzmap}
\end{figure}

\begin{figure}[!t]
\centering
\imgplaceholder{comparacion_lazo_abierto.png}
\caption{Comparaci\'on en lazo abierto entre el modelo lineal y el modelo realista.}
\label{fig:ol}
\end{figure}

\begin{figure}[!t]
\centering
\imgplaceholder{seguimiento_y_perturbacion.png}
\caption{Seguimiento de referencia y rechazo de perturbaciones en lazo cerrado.}
\label{fig:cl}
\end{figure}

\begin{figure}[!t]
\centering
\imgplaceholder{esfuerzo_de_control.png}
\caption{Se\~nal de control generada por el controlador IMC.}
\label{fig:u}
\end{figure}

\section{Implementaci\'on}
La carpeta entregable contiene scripts MATLAB para el an\'alisis, el dise\~no del controlador, la simulaci\'on lineal y no lineal, adem\'as de un generador autom\'atico de modelos Simulink:
\begin{itemize}
\item \texttt{Modelado\_Nivel\_Esclusa.m},
\item \texttt{parametros\_esclusa.m},
\item \texttt{analisis\_nivel\_esclusa.m},
\item \texttt{generar\_modelos\_simulink\_esclusa.m},
\item \texttt{run\_all.m}.
\end{itemize}

\section{Conclusiones}
El sistema propuesto satisface el requisito de una planta SISO de tercer orden y permite estudiar de manera coherente la diferencia entre un modelo lineal de dise\~no y un modelo realista con no linealidades hidr\'aulicas. El dise\~no IMC entrega una estructura simple para control de nivel y deja una base utilizable para ajustes posteriores en MATLAB y Simulink.

\end{document}
